% from main.tex
%%%%%%%%%%  scan idx 156  |  folio 143  %%%%%%%%%%
abélien, et l'on a un homomorphisme
\[
  \rho:\Cart_S(\mathcal C^{\is},\underline F)
  \longrightarrow
  \Cart_S(\Tr(\mathcal C)^{\is},\underline F),
\]
défini par \(\rho(f)(T)=f(L)-f(L')-f(L'')\) pour \(f\in\Cart_S(\mathcal C^{\is},\underline F)\) et \(T=(L'\to L\to L''\to L'[1])\in\Ob\Tr(\mathcal C)\). Dire que \(f\in\Cart_S(\mathcal C^{\is},\underline F)\) est additif signifie que \(\rho(f)=0\), autrement dit on a une suite exacte (canonique fonctorielle)
\begin{equation}\tag{6.6.1}
  0\longrightarrow \Add_S(\mathcal C,F)
  \longrightarrow \Cart_S(\mathcal C^{\is},\underline F)
  \xrightarrow{\rho}
  \Cart_S(\Tr(\mathcal C)^{\is},\underline F).
\end{equation}

\textbf{b)} Soit \(u:X\to Y\) une flèche de \(S\). Les foncteurs « image inverse par \(u\) », étant tous isomorphes entre eux, définissent (6.3) un unique homomorphisme \(k(Y)\longrightarrow k(X)\). Par suite, \(X\mapsto k(X)\) définit un préfaisceau abélien sur \(S\), noté \(k(\mathcal C)\). Une \(S\)-fonction additive de \(\mathcal C\) dans \(F\) s'interprète alors simplement comme un homomorphisme de préfaisceaux \(k(\mathcal C)\longrightarrow F\), autrement dit on a
\[
  \Add_S(\mathcal C,F)=\Hom(k(\mathcal C),F).
\]
Le préfaisceau \(k(\mathcal C)\) dépend bien entendu fonctoriellement de \(\mathcal C\) vis-à-vis des \(S\)-foncteurs additifs (resp. exacts).

\textbf{Proposition 6.7.} Soient \(S\) un site, \(\mathcal C\) une \(S\)-catégorie abélienne plate, \(\mathcal C_0\) une sous-\(S\)-catégorie additive strictement pleine de \(\mathcal C\) satisfaisant aux hypothèses (4.0). Alors le morphisme de préfaisceaux
\[
  k(\mathcal C_0)\longrightarrow k(\Parf(\mathcal C,\mathcal C_0))\qquad(\text{cf. }(4.9)),
\]
induit par l'inclusion de \(\mathcal C_0\) dans \(\Parf(\mathcal C,\mathcal C_0)\), définit un isomorphisme sur les faisceaux associés. En d'autres termes, pour tout faisceau abélien \(F\) sur \(S\), l'%
